\documentclass[12pt]{article}
\usepackage{amsmath,amssymb,amsthm}
\usepackage{geometry}
\geometry{margin=1in}
\usepackage{hyperref}
\usepackage{booktabs}
\usepackage{cite}
\usepackage{authblk}

\newtheorem{theorem}{Theorem}
\newtheorem{corollary}[theorem]{Corollary}
\newtheorem{proposition}[theorem]{Proposition}
\newtheorem{lemma}[theorem]{Lemma}
\theoremstyle{definition}
\newtheorem{definition}[theorem]{Definition}
\newtheorem{remark}[theorem]{Remark}
\newtheorem{conjecture}[theorem]{Conjecture}

\newcommand{\Z}{\mathbb{Z}}
\newcommand{\Q}{\mathbb{Q}}
\newcommand{\R}{\mathbb{R}}
\newcommand{\C}{\mathbb{C}}
\newcommand{\ph}{\varphi}
\newcommand{\Lk}{L_k}
\newcommand{\Lcal}{\mathcal{L}}
\newcommand{\logph}{\log\ph}
\newcommand{\tr}{\mathrm{tr}}
\newcommand{\ITF}{\mathrm{ITF}}

\title{Lucas Numbers in the Geodesic Length Spectrum and\\
Covering Tower Homology of a Compact Hyperbolic 3-Manifold\\
\large (Version 3: Conditional Field Obstruction and Citation
Corrections)}
\author[1]{Marvin L.\ Gentry}
\affil[1]{Independent Researcher, Seattle, WA\\
\texttt{drmlgentry@protonmail.com} \quad ORCID: 0009-0006-4550-2663}
\date{\today}

\begin{document}
\maketitle

\begin{abstract}
We give a corrected and strengthened account of how the Lucas
sequence $L_k$ and the golden ratio $\ph=(1+\sqrt5)/2$ relate to the
geometry and covering-tower arithmetic of
$M=\texttt{m003}(-2,3)$ (\texttt{OrientableClosedCensus[1]},
$\mathrm{vol}\approx0.9814$, $H_1\cong\Z/5$, $\ITF(M)=\Q(\sqrt{-3})$).

\textbf{Binet bifurcation} (Lemma~\ref{lem:bifurcation}). For every
integer $k\geq0$,
$\ph^k+\ph^{-k}=2\cosh(k\logph)$ equals the integer Lucas number $L_k$
if $k$ is even, and equals $\sqrt5\,F_k$ (Fibonacci) if $k$ is odd.
This single fact was previously mis-stated as ``$=L_k$ for all $k$,''
an error traced to conflating the recurrence-defined integer sequence
$L_k$ with the continuous function $\ph^k+\ph^{-k}$.

\textbf{Algebraic trace bridge} (Theorem~\ref{thm:bridge}), corrected.
If $\gamma\in\pi_1(M)$ has holonomy with real eigenvalues $\ph^{\pm k}$
($k\in\Z_{\geq1}$), then $\tr(\rho(\gamma))=\pm L_k$ for even $k$ and
$\pm\sqrt5 F_k$ for odd $k$.

\textbf{Field obstruction} (Corollary~\ref{cor:obstruction}), conditional
on Hypothesis~H ($k(\Gamma)=\ITF(\Gamma)$; status open,
Remark~\ref{rem:hypothesis}). Since $\ITF(M)=\Q(\sqrt{-3})$
and $\Q(\sqrt{-3})\cap\Q(\sqrt5)=\Q$, \emph{no} $\gamma\in\pi_1(M)$
can have real eigenvalues $\ph^{\pm k}$ for odd $k\geq1$, \emph{provided}
$k(\Gamma)=\ITF(\Gamma)$. The
``odd rungs'' of the golden-ratio length ladder
($\ell=2k\logph$, $k$ odd) are forbidden under this hypothesis.

\textbf{The covering tower: one Lucas prime}
(Theorems~\ref{thm:cover}--\ref{thm:why11}), corrected. Through
degree~6, the algebraic covering tower of $M$
(via \texttt{M.covers(d)}) has exactly one cover: a degree-5 cover
with $H_1\cong(\Z/11)^2$. The cover prime $11=L_5$ is explained by the
Farey cover-prime formula $p_n=5n(n+1)+1$ and the congruence
$p_n\equiv1\pmod5$, which forces $11=L_5$ to appear precisely at
$n=1$. The previously claimed ``prime dictionary''
$\{2,3,7,11,29\}=\{L_0,L_2,L_4,L_5,L_7\}$ conflated this algebraic
cover with census-volume commensurability coincidences.
\end{abstract}

\section{Introduction}

The arithmetic topology program studied in
\cite{gentry-ckm,gentry-pmns,gentry-cp,gentry-twist}
proposes that Standard Model flavor parameters arise from the geometry
of two compact hyperbolic 3-manifolds:
\begin{align*}
M_{\mathrm{PMNS}} &= \texttt{OrientableClosedCensus[1]},\quad
  \mathrm{vol} = 0.98137,\quad H_1 \cong \mathbb{Z}/5,\\
M_{\mathrm{CKM}}  &= \texttt{OrientableClosedCensus[43]},\quad
  \mathrm{vol} = 2.02885,\quad H_1 \cong \mathbb{Z}/5.
\end{align*}
Both are Dehn fillings of the two smallest-volume cusped orientable
hyperbolic 3-manifolds $m003(-2,3)$ and $m006(-5,2)$. Their
holonomy representations $\rho\colon \pi_1(M) \to \mathrm{PSL}(2,\C)$
have invariant trace field $\Q(\sqrt{-3})$~\cite{maclachlan-reid}.

The fermion mass lattice~\cite{gentry-core} postulates that Standard
Model fermion masses satisfy $m = m_e \cdot \ph^{q/4}$ for $q \in \Z$.
The base $\ph$ is independently derived from the Mahler measure identity
$\log M(\Delta_{m003}) = 2\logph$~\cite{gentry-jktr}.

This raises a question: the trace field is $\Q(\sqrt{-3})$, yet the
mass lattice involves $\ph \in \Q(\sqrt{5})$. How, if at all, does
$\Q(\sqrt{5})$ enter the arithmetic of a $\Q(\sqrt{-3})$ manifold?
This paper's earlier version (v1) answered ``via the Lucas sequence,
exactly, for all $k$.'' That answer rested on an identity that holds
only for even $k$. The corrected answer is more interesting: for odd
$k$, the would-be $\sqrt5$-valued trace is \emph{excluded}, and this
exclusion is itself the resolution.

\section{The Binet Bifurcation and the Trace Bridge}

\begin{definition}[Lucas and Fibonacci sequences]
The \emph{Lucas sequence} satisfies $L_0=2$, $L_1=1$,
$L_k=L_{k-1}+L_{k-2}$; the \emph{Fibonacci sequence} satisfies
$F_0=0$, $F_1=1$, $F_k=F_{k-1}+F_{k-2}$. Equivalently, with
$\psi=(1-\sqrt5)/2=-1/\ph$,
\[
  L_k=\ph^k+\psi^k, \qquad F_k=\frac{\ph^k-\psi^k}{\sqrt5}.
\]
The first values are $L: 2,1,3,4,7,11,18,29,\ldots$ and
$F: 0,1,1,2,3,5,8,13,\ldots$.
\end{definition}

\begin{lemma}[Binet bifurcation]
\label{lem:bifurcation}
For every $k\in\Z_{\geq0}$,
\[
  \ph^k+\ph^{-k} \;=\; 2\cosh(k\logph) \;=\;
  \begin{cases}
    L_k, & k \text{ even},\\[2pt]
    \sqrt5\,F_k, & k \text{ odd}.
  \end{cases}
\]
\end{lemma}

\begin{proof}
Since $\psi=-1/\ph$, $\psi^k=(-1)^k\ph^{-k}$. If $k$ is even,
$\psi^k=\ph^{-k}$, so $L_k=\ph^k+\psi^k=\ph^k+\ph^{-k}$ directly. If
$k$ is odd, $\psi^k=-\ph^{-k}$, so $\ph^k-\psi^k=\ph^k+\ph^{-k}$, and
dividing by $\sqrt5$ in the definition of $F_k$ gives
$\sqrt5F_k=\ph^k-\psi^k=\ph^k+\ph^{-k}$.
\end{proof}

\noindent\emph{Numerical check} (to 50 digits, all differences
$<10^{-48}$): $k=1$: $\ph+\ph^{-1}=\sqrt5=\sqrt5\cdot F_1$. $k=3$:
$\ph^3+\ph^{-3}=2\sqrt5=\sqrt5\cdot F_3$ ($F_3=2$). $k=5$:
$\ph^5+\ph^{-5}=5\sqrt5=\sqrt5\cdot F_5$ ($F_5=5$). Note in particular
$\ph^5+\ph^{-5}\approx11.180\neq11=L_5$.

\begin{theorem}[Algebraic trace bridge, corrected]
\label{thm:bridge}
Let $\gamma\in\pi_1(M)$ for $M$ a hyperbolic 3-manifold with holonomy
$\rho\colon\pi_1(M)\to\mathrm{PSL}(2,\C)$. Suppose $\rho(\gamma)$ has
real eigenvalues $\ph^{\pm k}$ for some $k\in\Z_{\geq1}$ (equivalently,
$\gamma$ is loxodromic with vanishing rotational part, complex length
$2k\logph$). Then
\[
  \tr(\rho(\gamma)) = \pm(\ph^k+\ph^{-k}) =
  \begin{cases}
    \pm L_k \in \Z, & k \text{ even},\\[2pt]
    \pm\sqrt5\,F_k \in \sqrt5\,\Z, & k \text{ odd}.
  \end{cases}
\]
\end{theorem}

\begin{proof}
For $A\in\mathrm{SL}(2,\C)$ with eigenvalues $\lambda,\lambda^{-1}$,
$\tr(A)=\lambda+\lambda^{-1}$. Take $\lambda=\ph^k$ and apply
Lemma~\ref{lem:bifurcation}; the sign is the usual
$\mathrm{SL}(2,\C)\to\mathrm{PSL}(2,\C)$ lift ambiguity.
\end{proof}

\begin{corollary}[Field obstruction --- the even-Lucas selection rule,
conditional]
\label{cor:obstruction}
Let $M=M_{\mathrm{PMNS}}$, with $\ITF(\Gamma)=\Q(\sqrt{-3})$
\cite{maclachlan-reid}. Suppose in addition that
\[
  \textbf{Hypothesis H:} \qquad k(\Gamma) = \ITF(\Gamma),
\]
i.e.\ $\Gamma=\pi_1(M)$ lifts to $\mathrm{SL}(2,\Q(\sqrt{-3}))$ without
a quadratic extension. Then no $\gamma\in\pi_1(M)$ satisfies the
hypothesis of Theorem~\ref{thm:bridge} for odd $k\geq1$.
\end{corollary}

\begin{proof}
Under Hypothesis H, $\tr(\rho(\gamma))\in k(\Gamma)=\Q(\sqrt{-3})$ for
all $\gamma\in\pi_1(M)$ (for an appropriate lift of $\rho$). If
$\gamma$ satisfied the hypothesis of Theorem~\ref{thm:bridge} for odd
$k$, then $\tr(\rho(\gamma))=\pm\sqrt5F_k$. For $k\geq1$, $F_k\geq1$,
so $\pm\sqrt5F_k\notin\Q$. Since $\Q(\sqrt{-3})$ and $\Q(\sqrt5)$ are
distinct quadratic fields, $\Q(\sqrt{-3})\cap\Q(\sqrt5)=\Q$, so
$\pm\sqrt5F_k\notin\Q(\sqrt{-3})$, contradicting Hypothesis H.
\end{proof}

\begin{remark}[Status of Hypothesis H --- the paper's primary open item]
\label{rem:hypothesis}
$\ITF(\Gamma)=\Q(\sqrt{-3})$ is a \emph{commensurability invariant} and
is well established for $M_{\mathrm{PMNS}}$ via its relation to the
Eisenstein/figure-eight arithmetic~\cite{maclachlan-reid}. The
(non-invariant) trace field $k(\Gamma)=\Q(\tr\rho(\gamma):\gamma\in\Gamma)$
is \emph{not} a commensurability invariant, and in general
$[k(\Gamma):\ITF(\Gamma)]\in\{1,2\}$: $k(\Gamma)$ may equal $\ITF(\Gamma)$,
or be a quadratic extension $\ITF(\Gamma)(\sqrt{d})$ for some
$d\in\ITF(\Gamma)$. Hypothesis H asserts the former for the specific
group $\Gamma=\pi_1(M_{\mathrm{PMNS}})$.

A 10-digit numerical holonomy computation found
$\tr(\rho(aa))\approx-2.7881$. Any \emph{real} element of $\Z[\omega]$
(the ring of integers of $\Q(\sqrt{-3})$) is a rational integer, since
$\mathrm{Im}(a+b\omega)=b\sqrt3/2=0\iff b=0$; $-2.7881$ is not an
integer. This is \emph{suggestive} that $k(\Gamma)\supsetneq\Q(\sqrt{-3})$
(degree 4), though not conclusive at 10 digits, and `aa' need not be
representative of the field-generating set.

If Hypothesis H fails, Corollary~\ref{cor:obstruction}'s conclusion
depends on a further question: does the resulting quartic
$k(\Gamma)=\Q(\sqrt{-3})(\sqrt{d})$ contain $\sqrt5$? If
$d\equiv5\pmod{(\Q(\sqrt{-3})^\times)^2}$, it does, and the corollary
fails outright. If not,
$\Q(\sqrt{-3})(\sqrt{d})\cap\Q(\sqrt5)=\Q$
still (two non-isomorphic biquadratic fields sharing only $\Q$ unless
$d$ is chosen to force it), and an analogue of the corollary may be
recoverable with $\ITF(\Gamma)$ replaced by $k(\Gamma)$ throughout.

We regard Hypothesis H as the single open item on which
Corollary~\ref{cor:obstruction} -- and hence the central claim of this
paper -- currently rests. A high-precision ($\geq30$ digit)
recomputation of several generator traces, followed by
\texttt{algdep}, would settle it: either $k(\Gamma)=\Q(\sqrt{-3})$
(Corollary~\ref{cor:obstruction} stands as proved, unconditionally),
or $k(\Gamma)=\Q(\sqrt{-3},\sqrt{d})$ for an explicit $d$ (the
corollary stands, fails, or generalizes, according to whether
$5d\in(\Q(\sqrt{-3})^\times)^2$).
\end{remark}

\begin{remark}[What this resolves, and what it does not]
\label{rem:resolution}
One part of this is unconditional: $\sqrt5\notin\ITF(\Gamma)=\Q(\sqrt{-3})$
is pure field theory, independent of Hypothesis H. The even-$k$ rungs
($\ell=2k\logph$, $k$ even) would require only $L_k\in\Z\subset\Q(\sqrt{-3})$
--- no obstruction at any level, but also no existence claim; whether
such $\gamma$ actually occur in $\pi_1(M)$ is a separate, open,
empirical question. Whether the odd-$k$ rungs are \emph{actually}
absent is exactly Hypothesis H's content (Remark~\ref{rem:hypothesis}):
\emph{if} $k(\Gamma)=\ITF(\Gamma)$, they are provably absent by
Corollary~\ref{cor:obstruction}; if $k(\Gamma)\supsetneq\Q(\sqrt{-3})$,
their fate depends on whether the extension contains $\sqrt5$. Either
way this is a \emph{selection rule on a discrete length spectrum,
falsifiable in principle}: by Theorem~\ref{thm:bridge} (unconditional),
any $\gamma\in\pi_1(M_{\mathrm{PMNS}})$ with real eigenvalues
$\ph^{\pm k}$ for odd $k$ has $\tr(\rho(\gamma))=\pm\sqrt5F_k\notin\Q(\sqrt{-3})$
--- and since $\ITF(\Gamma)=\Q(\sqrt{-3})$ is independently established,
finding such a $\gamma$ would directly refute Hypothesis H (show
$k(\Gamma)\neq\ITF(\Gamma)$), not Theorem~\ref{thm:bridge} or the trace
field identification.
\end{remark}

\begin{remark}[Where does $L_5=11$ fit?]
\label{rem:L5}
$L_5=11$ is odd-indexed. By Corollary~\ref{cor:obstruction} (conditional
on Hypothesis H, Remark~\ref{rem:hypothesis}), $11$
cannot be $|\tr(\rho(\gamma))|$ for $\gamma$ with eigenvalues
$\ph^{\pm5}$ (length $10\logph$). If $11$ appears as a trace value at
all --- and as a rational integer there is no field obstruction to it
appearing \emph{somewhere} --- it does so at a geometrically generic
length, $\ell=2\,\mathrm{arccosh}(11/2)\approx4.7791$, with
$\ell/\logph\approx9.931\notin\Q$. Section~\ref{sec:cover} shows that
$L_5=11$'s role in this programme is entirely number-theoretic ---
as the algebraic cover prime of $M_{\mathrm{PMNS}}$ --- and is
unrelated to the geodesic-length bridge above. These are two
independent phenomena that happen to share the symbol $11=L_5$.
\end{remark}

\begin{remark}[Geodesic spectrum of $M_{\mathrm{PMNS}}$]
\label{rem:spectrum}
For $M_{\mathrm{PMNS}}=\texttt{m003}(-2,3)$, the first ten geodesic
lengths $\ell$ (real parts) give $\ell/(2\logph)\in
\{0.60,0.75,0.92,1.04,1.08,1.86,1.89,2.19,2.31,2.35\}$ --- none an
integer. This is consistent with (does not contradict)
Theorem~\ref{thm:bridge}: exact integer-$k$ real-eigenvalue elements
are a special, evidently rare, condition, not a generic feature of
the spectrum. The entry near $0.75$ corresponds to $k=3/4$, discussed
next.
\end{remark}

\begin{remark}[Quarter-integer values: numerical, not algebraic-integer]
\label{rem:quarter}
For real $x\geq0$, define $\Lcal(x):=\ph^x+\ph^{-x}=2\cosh(x\logph)$.
For $x\in\Z$ this is Lemma~\ref{lem:bifurcation}'s bifurcated
quantity. For $x\notin\Z$, $\Lcal(x)$ is an algebraic number in
$\Q(\ph^{1/4})$ (degree $\leq8$ over $\Q$) but --- being a real field
that does not contain $\sqrt{-3}$ --- $\Q(\ph^{1/4})\cap\Q(\sqrt{-3})
=\Q$. By the same structure as Corollary~\ref{cor:obstruction}, this
gives $\Lcal(x)\notin\ITF(\Gamma)$ for irrational $x$ \emph{provided}
$k(\Gamma)=\ITF(\Gamma)$ (Hypothesis H, Remark~\ref{rem:hypothesis});
the unconditional fact is only $\Lcal(x)\notin\Q$ for $x\in(0,1)$.
Independently of Hypothesis H, the previously-noted match
$\sigma_{\mathrm{opt}}\approx\tfrac32\logph$ (i.e.\ $x=3/4$, agreement
$\approx0.004\%$, \emph{not} $0$) is empirically a numerical
coincidence rather than an exact identity regardless --- a $0.004\%$
residual is itself evidence against
$|\tr(\gamma_2)|=\Lcal(3/4)$ holding exactly, independent of which
field $\Lcal(3/4)$ turns out to live in. We use $\Lcal$ (not $L$) for
non-integer arguments throughout to avoid the notational collision
that caused the v1 error.
\end{remark}

\section{The Covering Tower: $11=L_5$}
\label{sec:cover}

\begin{theorem}[The unique algebraic cover prime, corrected]
\label{thm:cover}
Through degree~6, \texttt{M.covers(d)} applied to $M_{\mathrm{PMNS}}$
returns exactly one cover: a degree-5 cover with
$H_1\cong(\Z/11)^2$. Equivalently, $11=L_5$ is the unique new
algebraic cover prime of $M_{\mathrm{PMNS}}$ through degree~6.
\end{theorem}

\begin{proof}
Direct computation via \texttt{SnapPy}'s \texttt{M.covers(d)} for
$d=2,\ldots,6$~\cite{snappy}.
\end{proof}

\begin{remark}[Correction to the v1 ``prime dictionary'']
\label{rem:dict-correction}
Version 1 of this paper claimed
$\{2,3,7,11,29\}=\{L_0,L_2,L_4,L_5,L_7\}$ as a ``prime dictionary'' of
the covering tower. This conflated the algebraic cover prime $11$
(Theorem~\ref{thm:cover}) with primes $\{2,3,7,29\}$ that appear
because certain census manifolds at integer multiples of
$\mathrm{vol}(M_{\mathrm{PMNS}})$ are \emph{commensurable neighbours}
of $M_{\mathrm{PMNS}}$ --- not covers of it. Only $11=L_5$ is an
algebraic cover prime.
\end{remark}

\begin{theorem}[Why $11=L_5$, and only $11$]
\label{thm:why11}
Let $M_n=\texttt{m003}(-(n+1),2n+1)$ for $n\geq1$ be the Farey tower
containing $M_{\mathrm{PMNS}}=M_1$. Each $M_n$ has $H_1\cong\Z/5$ and
a unique degree-5 cover with $H_1\cong(\Z/p_n)^2$, where
\[
  p_n = 5n(n+1)+1.
\]
Then $p_n\equiv1\pmod5$ for all $n$, and $11=L_5$ divides $p_n$ if and
only if $n\equiv1$ or $9\pmod{11}$. In particular $p_1=11=L_5$
exactly (the $n=1$ case of the divisibility, where $p_1$ \emph{is}
the Lucas prime), while $p_2=31,\ldots,p_{10}=541$ are not Lucas
numbers. $11=L_5$ is thus the \emph{unique} Lucas prime in the
sequence $(p_n)$, and $M_{\mathrm{PMNS}}=M_1$ is the unique element
of its own Farey tower with this property.
\end{theorem}

\begin{proof}
See the companion paper~\cite{gentry-z5bridge}, Theorems 2.5, 5.2,
and 5.3: the cover-prime formula, the forbidden-Lucas-primes result
(ruling out $L_0,L_2,L_4\in\{2,3,7\}$ for all $n$ by residues mod
2,3,7), and the first-permitted-Lucas-prime result
($11=L_5\mid p_n\iff n\equiv1,9\bmod11$).
\end{proof}

\begin{remark}[On the degree-9 scan and ``Lucas-purity'']
\label{rem:purity-status}
Version 1 additionally reported a degree-9 scan (24 covers, 47
torsion computations) finding only the Lucas-prime torsion set
$\{2,3,7,11,29\}$, and conjectured this as ``Lucas-purity.'' That
scan predates the distinction drawn in
Remark~\ref{rem:dict-correction} between algebraic covers and
commensurability-class coincidences, and has not been independently
re-examined under that distinction at degrees~7--9. We do not
reassert the Lucas-purity conjecture here. In its place,
Theorem~\ref{thm:why11} gives a narrower but \emph{proved} statement:
within the Farey tower $(M_n)_{n\geq1}$, exactly one cover prime
($p_1=11=L_5$) is a Lucas number, and this is forced by elementary
congruences, not a purity phenomenon requiring further conjecture.
A re-scan of degrees~7--9 under the corrected algebraic/commensurable
distinction remains a worthwhile open computation.
\end{remark}

\begin{remark}[Quantitative control computation]
\label{rem:controls}
Independent of the corrections above, the following control
computation remains valid as a measurement of the rarity of
low-torsion-prime covers: a scan of 73 fillings across 5 cusped
manifolds (degree~$\leq6$) found non-Lucas torsion primes
$13,17,19,23,31$ appearing in $17$--$93\%$ of control fillings,
versus $0\%$ for $M_{\mathrm{PMNS}}$ at its single cover. (Through
degree~6, \texttt{M.covers(d)} returns no covers at all for
$M_{\mathrm{CKM}}$, so its corresponding figure is $0$ of $0$ ---
vacuously consistent, not a data point.) With at most one cover per
flavor manifold through degree~6, this comparison has limited
statistical power and should be read as suggestive rather than
conclusive.
\end{remark}

\section{Discussion}

The corrected picture is, if anything, more structured than the v1
claim. Rather than ``Lucas numbers appear everywhere, exactly,'' we
have: (1) a \emph{candidate forbidden} set of geodesic lengths
(Corollary~\ref{cor:obstruction}, conditional on Hypothesis H,
Remark~\ref{rem:hypothesis}), arising from a clean field-theoretic
obstruction and directly linking this paper to the Eisenstein-field
result $\ITF(M_{\mathrm{PMNS}})=\Q(\sqrt{-3})$; and (2) a single,
arithmetically-forced cover prime $11=L_5$
(Theorem~\ref{thm:why11}), explained by elementary congruences on the
Farey cover-prime sequence rather than by an open-ended purity
conjecture. Item (2) is unconditional; item (1) is the paper's main
claim and rests on Hypothesis H.

\paragraph{TQFT remark (corrected).} The identity
$\Lcal(k)=2\cosh(k\logph)$ relates $\Lcal$ to Chebyshev polynomials of
the first kind, equal to $L_k$ (even $k$) or $\sqrt5F_k$ (odd $k$) by
Lemma~\ref{lem:bifurcation}. In $\mathrm{SU}(2)$ Chern--Simons theory
at level $r=3$ (Fibonacci anyons), $q=e^{2\pi i/5}$ gives
$[2]_q=2\cos(2\pi/5)=\ph-1$, a trigonometric, not hyperbolic, cosine;
the analogy remains formal.

\paragraph{Open directions.} (0) Settling Hypothesis H itself
(Remark~\ref{rem:hypothesis}) via high-precision \texttt{algdep} ---
the prerequisite for the paper's central claim. (1) Does
$\pi_1(M_{\mathrm{PMNS}})$ contain
\emph{any} element with real eigenvalues $\ph^{\pm k}$ for
even $k\geq2$ (the rungs Corollary~\ref{cor:obstruction} permits but
does not require)? (2) A re-scan of algebraic covers at
degrees~7--9 under the corrected algebraic/commensurable distinction.
(3) The dynamical origin of the mass lattice $m=m_e\ph^{q/4}$, now
understood as a statement about $\Lcal(q/8)\in\Q(\ph^{1/4})$
(Remark~\ref{rem:quarter}) rather than about integer Lucas traces.

\section*{Data Availability}
All computations use SnapPy~\cite{snappy} v3.3.2.\\
Scripts: \url{https://github.com/drmlgentry/hyperbolic-flavor-scan}. The Farey
cover-prime results of Theorem~\ref{thm:why11} are proved in
\cite{gentry-z5bridge}.

\section*{Funding}
This research received no external funding.

\section*{Conflict of Interest}
The author declares no conflicts of interest.

\section*{Generative AI Disclosure}
Claude was used for \LaTeX{} formatting assistance, code debugging,
identification of the bifurcation error corrected in v2, and
identification of the trace-field gap (Hypothesis H,
Remark~\ref{rem:hypothesis}) addressed in v3. All mathematical
results were conceived, proved, and verified by the author.

\begin{thebibliography}{99}

\bibitem{gentry-ckm}
M.~L.~Gentry,
\textit{Quark Mixing from Hyperbolic Geometry},
Results in Physics, submitted (2026). SSRN:10.2139/ssrn.6583550.

\bibitem{gentry-pmns}
M.~L.~Gentry,
\textit{Lepton Mixing from the Borel Structure of Hyperbolic Holonomy},
Results in Physics, submitted (2026). SSRN:10.2139/ssrn.6583553.

\bibitem{gentry-cp}
M.~L.~Gentry,
\textit{CP Violation from A-Factor Twist Angles},
Results in Physics, submitted (2026).

\bibitem{gentry-twist}
M.~L.~Gentry,
\textit{Twist Angle Spectrum of Hyperbolic Holonomy},
Nuclear Physics B, submitted (2026).

\bibitem{gentry-core}
M.~L.~Gentry,
\textit{Hyperbolic Flavor Geometry: A Unified Framework},
preprint (2026).

\bibitem{gentry-jktr}
M.~L.~Gentry,
\textit{The Alexander Polynomial of the $(-2,3,7)$ Pretzel Knot
and the Golden Ratio}, JKTR, submitted (2026).
SSRN:10.2139/ssrn.6583549.

\bibitem{gentry-z5bridge}
M.~L.~Gentry,
\textit{The $\Z/5$ Bridge: Hyperbolic Covering Towers, Rational
Torsion of $X_0(11)$, and an Arithmetic Proof of the Universal Cover
Prime Formula}, preprint (2026).

\bibitem{maclachlan-reid}
C.~Maclachlan and A.~W.~Reid,
\textit{The Arithmetic of Hyperbolic 3-Manifolds},
Springer (2003).

\bibitem{snappy}
M.~Culler, N.~M.~Dunfield, M.~Goerner, J.~R.~Weeks,
\textit{SnapPy}, \url{http://snappy.computop.org}.

\end{thebibliography}

\end{document}
